% [原文 p.241]
\begin{theorem*}[{[Kol89]}]{6.15}
三维终端翻模保持解析奇点。确切地说，若 $X$ 是只具有终端奇点 $\{x_1,\dots,x_n\}$ 的三维簇，则其翻模 $X'$ 也具有同样多的奇点 $\{x'_1,\dots,x'_n\}$，并且（适当重新编号后）对每个 $i$，作为解析芽有同构 $(X,x_i)\cong(X',x'_i)$。
\end{theorem*}

% [Tutor Note: 定理 6.15 说明翻模不改变奇点的解析类型；正因如此，三维终端奇点的局部模型在翻模下不变，标准翻模的构造才能逐点归约到这些局部模型上，这也是标准翻模可终止的关键。]

接下来处理翻模的终止性。为使陈述明确，考虑如下设定。

\begin{definition}{6.16}
所谓 $(K_X+\Delta)$-翻转的列（sequence of $(K_X+\Delta)$-flips），是指对 $(X^i,\Delta^i)$ 与双有理映射 $\varphi^i:X^i\dashrightarrow X^{i+1}$ 的列：它从对 $(X^1,\Delta^1)=(X,\Delta)$ 开始，其中 $\varphi^i:X^i\dashrightarrow X^{i+1}$ 是 $(K_{X^i}+\Delta^i)$-翻转，且 $\Delta^{i+1}:=\varphi^i_*\Delta^i$ 依次确定。若 $\varphi^i$ 对所有 $i\ge1$ 都有定义，则称之为翻转的无限列。
\end{definition}

\begin{theorem*}[三维标准翻转的终止]{6.17}
三维标准 $K_X$-翻转的任意列都是有限的。此外，三维射线型标准 $(K_X+\Delta)$-翻转的任意列都是有限的。
\end{theorem*}

\begin{remark}{6.18}
若 $X$ 是解析空间，则在定理 6.17 与推论 6.19 中须假设 $X$ 是紧致子集的充分小邻域。这是为了避免翻转的曲线不断远离任意紧致子集的情形。
\end{remark}

由于翻模是特殊的翻转，故得到如下结果。

\begin{corollary}[三维终端翻模的终止]{6.19}
三维射线型终端 $(K+D)$-翻模的任意列都是有限的。
\end{corollary}

定理 6.17 的证明使用关于如下不变量的归纳法。

\begin{definition}[{[Kol89, Kaw92b, K+92]}]{6.20}
设 $(X,\Delta=\sum a_iD_i)$ 是标准对，$D_i$ 是互不相同的素除子。置 $a:=\max\{a_i\}$、$S:=\sum a_i\Z_{\ge0}\subset\Q$（但若 $\Delta=0$，则置 $a=0$、$S=\{0\}$）。此时定义
\[
d(X,\Delta):=\sum_{\xi\in S,\ \xi\ge a}\#\{\text{$X$ 上的例外除子 }E \mid a(E,X,\Delta)<1-\xi\}.
\]
% [原文 p.242]
由于 $\operatorname{totaldiscrep}(X,\Delta)=-a$，由命题 2.36(2) 可知 $d(X,\Delta)<\infty$。这里还指出，$d(X,\Delta)$ 在作用 $(K_X+\Delta)$-翻转后不增加（引理 3.38）。
\end{definition}

\begin{lemmap}{6.21}
设 $\varphi:X\dashrightarrow X'$ 是三维标准对 $(X,\Delta=\sum_{i=1}^k a_iD_i)$ 的 $(K_X+\Delta)$-翻转。设 $C'\subset X'$ 是后翻转曲线（定义 3.33），$E_{C'}$ 是在 $C'$ 的泛点附近爆破 $C'$ 所得的例外除子。此时 $X'$ 沿 $C'$ 非奇异，且下式成立：
\[
0\le a(E_{C'},X,\Delta)<a(E_{C'},X',\Delta')=1-\sum a_i\cdot\operatorname{mult}_{C'}(D'_i).
\]
\end{lemmap}

\begin{proof}
$(X',\Delta')$ 沿 $C'$ 是终端的（引理 3.38）。因此 $X'$ 沿 $C'$ 非奇异。其余是显然的计算（参见引理 2.29）。
\end{proof}

\begin{proof}[定理 6.17 的证明]
记 $\Delta=\sum_{i=1}^k a_iD_i$，其中 $a_1\le\cdots\le a_k$。使用定义 6.16 的记号。对 $A\subset X$，以 $A^j$ 表示其在 $X^j$ 上的双有理变换。

证明对 $k$ 归纳。若 $k=0$，则由引理 3.38 与引理 6.21，对所有 $j$ 有 $d(X^{j-1},0)>d(X^j,0)$。于是 $k=0$ 的情形证毕。

设 $k>0$。由定义 6.20，$d(X^j,\Delta^j)$ 是 $j$ 的非减函数。若关于 $\varphi^{j-1}$ 的后翻转曲线 $C^j$ 含于 $D^j_k$ 中，则由引理 3.38 与引理 6.21 得 $a_k<1$ 且 $d(X^{j-1},\Delta^{j-1})>d(X^j,\Delta^j)$。因此当 $j\gg0$ 时，$D^j_k$ 不含后翻转曲线。以 $\bar D^j_k$ 表示 $D^j_k$ 的正规化，则 $\varphi^{j-1}$ 确定双有理映射 $\bar D^{j-1}_k\to\bar D^j_k$。固定任意 $j$，映射 $\bar D^j_k\to\bar D^l_k$（$l>j$）的例外曲线是数值独立的（引理 3.40）。因此对所有 $l\ge j\gg0$，有同构 $\bar D^j_k\cong\bar D^l_k$。也就是说，当 $j\gg0$ 时，$D^j_k$ 也不含前翻转曲线（定义 3.33）。若 $f^j:X^j\to Y^j$ 是射线型翻转收缩，则 Q-Cartier 除子 $D^j_k$ 在 $Y^j$ 上是 nef 的。于是 $-\bigl(K_{X^j}+\sum_{i=1}^{k-1}a_iD^j_i\bigr)$ 是 $f^j$-丰富的，且 $\varphi^j:X^j\dashrightarrow X^{j+1}$ 是射线型 $\bigl(K_{X^j}+\sum_{i=1}^{k-1}a_iD^j_i\bigr)$-翻转。这样便归结为 $k-1$ 的情形。
\end{proof}

% [原文 p.243]
\section{终端化与 Q-分解}
\takeaway{本节用射线型终端翻模证明两个定理：终端化定理把至多具有标准奇点的三维簇换成奇点更简单的终端模型，Q-分解定理则把 Weil 除子化为 Q-Cartier。}

本节使用前节研究的射线型终端翻模，证明关于终端化与 Q-分解的两个定理。

\begin{definition}[{[Rei83b]}]{6.22}
设 $(X,\Delta)$ 是对，$f:Y\to X$ 是双有理态射。若 $f$-例外的不可约除子 $E\subset Y$ 满足 $a(E,X,\Delta)=0$，则称 $E$（关于 $(X,\Delta)$）是 crepant 的。若所有 $f$-例外除子都是 crepant 的，则称 $f$（关于 $(X,\Delta)$）是 crepant 的。
\end{definition}

最初的主要结果是 [Rei83b] 的终端化定理。

\begin{theorem*}[终端化（terminalization）]{6.23}
设 $X$ 是只具有标准奇点的三维代数簇（或解析空间）。此时存在 crepant 双有理态射 $f:Y\to X$，其中 $Y$ 至多具有终端奇点，且 $f$ 是射影的（或 $Y$ 的任意紧致子集的适当邻域上的射影态射）。
\end{theorem*}

\begin{remark}{6.24}
以 $\pi_X(:=f):X^{ter}(:=Y)\to X$ 表示定理 6.23 的证明中构造的终端化，则它具有如下性质。
\begin{enumerate}[label=（\arabic*）,leftmargin=3.2em]
\item （与开嵌入的交换性）开嵌入 $g:X_1\to X_2$ 唯一地延拓为开嵌入 $g^{ter}:X_1^{ter}\to X_2^{ter}$，使 $g\circ\pi_{X_1}=\pi_{X_2}\circ g^{ter}$。
\item （与 $X\to X^{an}$ 的交换性）$(\pi_X)^{an}=\pi_{X^{an}}$。
\end{enumerate}
本书称 $X^{ter}$ 为基准终端化（standard terminalization）。
\end{remark}

% [Tutor Note: 页图作「基準端末化」，standard terminalization 的 standard 是「基准」而非「標準（canonical）」；本书「标准」一词专指 canonical 奇点类别，两者不同字亦不同义，切勿混同。]

第二个主要结果是 [Kaw88] 的 Q-分解定理。

\begin{theorem*}[Q-分解（Q-factorialization）]{6.25}
设 $X$ 是只具有终端奇点的三维代数簇（或解析空间）。此时存在双有理态射 $f:Y\to X$，其中 $Y$ 是终端且 Q-分解的（即 $Y$ 上整体定义的每个 Weil 除子都是 Q-Cartier 的）簇；$f$ 在余维 1 处是同构，并且是射影的（或在 $Y$ 的任意紧致子集的适当邻域上射影）。
\end{theorem*}

% [Tutor Note: “Q-分解”即英文的 Q-factorialization：把簇换成其上任一 Weil 除子都为 Q-Cartier 的模型。它与“Q-分解的”（Q-factorial）是同一概念的名词与形容词形式，勿与一般的“分解”混淆。]

\begin{remark}{6.26}
关于定理 6.25 所给出的 Q-分解 $f:Y\to X$，指出以下几点。若三维代数簇 $X$ 不是固有的，则相伴的解析空间 $Y^{an}$ 未必是 $X^{an}$ 的 Q-分解。同样，若 $U\subset X$ 是解析空间 $X$ 的开集，则 $f^{-1}(U)$ 未必是 $U$ 的 Q-分解（参见例 2.17）。
\end{remark}

% [原文 p.244]
定理 6.23 的证明需要如下结果。

\begin{theoremp}[{[Rei83b]}]{6.27}
设 $X$ 是只具有标准奇点的三维簇，且对每个奇点 $P\in X$，$(P\in X)$ 的指数 1 覆盖是 cDV 奇点。设 $C$ 是 $X$ 的奇点集的一维部分（赋予约化结构），$I$ 是 $X$ 中 $C$ 的定义理想。此时，
\begin{enumerate}[label=（\arabic*）,leftmargin=3.2em]
\item $R=R(X):=\bigoplus_{\nu\ge0}I^{(\nu)}$ 是正规的有限生成 $\cO_X$-代数层，且 $f:Y=\Proj_X R\to X$ 的任意纤维至多一维。
\item $Y$ 是标准的，$K_Y\equiv f^*(K_X)$，并且 $\dim f^{-1}(x)=1$（$\forall x\in C$）。
\end{enumerate}
\end{theoremp}

% [Tutor Note: $I^{(\nu)}$ 是符号幂而非普通幂：它只在 $C$ 的泛点处与 $I^{\nu}$ 一致，在 $C$ 的其余点处可以更大。定理 6.27 的证明正是靠“符号幂的有限生成性”把这个局部代数问题转成整体几何问题。]

这里，$I^{(\nu)}$ 是称为 $I$ 的 $\nu$ 次符号幂（$\nu$-th symbolic power）的凝聚层，是满足下述性质的唯一层：$I^{(\nu)}\supset I^{\nu}$，$\cO_X/I^{(\nu)}$ 的伴随点都是 $C$ 的泛点，而 $I^{(\nu)}/I^{\nu}$ 的任一伴随点都是 $C$ 的非泛点。

\begin{proof}
首先假定 (1) 并证明 (2)。

置 $R':=\bigoplus_{\nu}I^{\nu}$。对每个 $\nu$，$I^{(\nu)}/I^{\nu}$ 是长度有限的层，因为其支集是零维的。

由于 $R$ 有限生成，适当选取有限集 $S\subset X$，可使 $R|_{X\setminus S}=R'|_{X\setminus S}$。置 $f':Y':=\Proj_X R'\to X$，则 $f|_{X\setminus S}=f'|_{X\setminus S}$。可以把 $S$ 取得更大，使其包含 $X$ 的全部非 cDV 点（推论 5.40）。

任取 $x\in C\setminus S$，设 $x\in H\subset U_x\subset X$ 是一般超曲面截口，其中 $U_x$ 是使 $U_x\cap H\cap C=\{x\}$ 的适当开邻域。$x\in H$ 是 Du Val 奇点，且 $(f')^{-1}(H)=B_xH$。这里 $B_xH$ 正规，且 $K_{B_xH}=(f')^*K_H$ 成立（参见定理 4.20 的证明的 (1)）。由于 $f^{-1}(S)$ 的余维为 2，由此得到 $K_{Y'\setminus(f')^{-1}(S)}=(f')^*K_{X\setminus S}$ 以及 $K_Y\equiv f^*K_X$。因此由引理 2.30，$Y$ 是标准的。这样便由 (1) 导出了 (2)。

% [原文 p.245]
其次局部解析地证明 (1)。设 $(P\in X)$ 是解析芽。不妨假定 $P\in C$。此时只需对指数 1 覆盖证明 (1)，再取循环群商（参见命题 6.8），便证明了 $P\in C\subset X$ 情形的 (1)。特别地，不妨假定 $(P\in X)$ 是 cDV 奇点。适当选取到非奇异曲线的态射 $\pi:(P\in X)\to(0\in V)$，可使 $X_0:=\pi^{-1}(0)$ 只具有 Du Val 奇点，并且 $\dim X_0\cap C=0$。

其次注意，制作底曲线 $V$ 的有限 Galois 基变换 $V'\to V$ 后，只需证明 (1)。这是因为，若 $G$ 是 Galois 群，则 $X\times_V V'$ 仍是 cDV 的，且一旦 $R(X\times_V V')$ 的正规有限生成性得证，其 $G$-不变部分 $R(X)$ 的正规有限生成性也立刻可得。

于是不妨假定存在 $\pi$ 的射影同时奇点消解（定理 4.28）$h:Z\to X$，使对每个 $v\in V$，$Z_v\to X_v$ 是极小奇点消解。以 $F$ 表示由 $I\cO_Z$ 定义的 $Z$ 的闭子空间的除子部分。由引理 6.28，存在有效除子 $G$ 使 $G\sim-F$。注意：$G$ 在 $V$ 的泛点 $\xi$ 上是 $h$-nef 的，$\dim h^{-1}(P)\le1$，且 $K_Z=h^*(K_X)$。

由于 $Z$ 非奇异，对充分小的 $\varepsilon>0$，$(Z,\varepsilon G)$ 是终端的。运行 $X$ 上的 $(K_Z+\varepsilon G)$-MMP。由于 $K_Z\sim0$，MMP 过程中标准除子是显然的。又由 $\dim h^{-1}(P)\le1$ 以及 $G$ 在 $X\setminus\{P\}$ 上是 nef 的，MMP 过程中不出现除子收缩，只出现 $G$-翻转收缩。因此，由射线型终端翻模的存在性与终止性，得到 $h':Z'\to X$。此时 $G$ 的双有理变换 $G'$ 在 $Z'$ 上是 $h'$-nef 的，且 $h'$ 在余维 1 处是同构。由于 $G'-K_{Z'}$ 在 $X$ 上是 nef 且大的，由相对无基点性定理 3.24，对 $\nu\gg0$，$\nu G'$ 是 $h'$-自由的。特别地，$\bigoplus_{\nu\ge0}h'_*\cO_{Z'}(\nu G')$ 有限生成。正规性由如下一般事实得到：若 $Z'$ 是任意正规簇、$G'$ 是 $Z'$ 上的 Weil 除子，则 $\bigoplus_{\nu\ge0}H^0(Z',\cO_{Z'}(\nu G'))$ 是正规环（即使不知它是否为 Noether 环）；整体情形参见 [Har77, Ex. I.5.14]。

由前面在 (1)$\Rightarrow$(2) 的证明中所用的论证与推论 4.14，在 $X\setminus\{P\}$ 上

% [原文 p.246]
$h'_*\cO_{Z'}(\nu G')=I^{(\nu)}(=I^{\nu})$ 成立是已知的。任取点 $s\in\cO_{P,X}$。由 $\dim(h')^{-1}(P)\le1$，有 $s\in h'_*\cO_{Z'}(\nu G')\iff s|_{X\setminus\{P\}}\in h'_*\cO_{Z'}(\nu G')|_{X\setminus\{P\}}$。又，准素理想 $I^{(\nu)}$ 也具有同样的性质。由此得 $h'_*\cO_{Z'}(\nu G')=I^{(\nu)}$，从而证明了 (1) 给出的 $R$ 正规且有限生成。由于 $h'$ 的纤维维数至多 1，(1) 给出的 $f$ 也具有同样的性质。
\end{proof}

\begin{lemmap}{6.28}
设 $f:U\to V$ 是固有双有理态射，$D$ 是 $U$ 上的 Cartier 除子。假设 $V$ 在代数情形是拟射影的、在解析情形是 Stein 的。此时存在 $U$ 上的有效 Cartier 除子 $F$ 使 $F\sim_f D$。
\end{lemmap}

\begin{proof}
设 $H$ 是 $V$ 上的 Cartier 除子，使 $f_*\cO_U(D)\otimes\cO_V(H)$ 具有非零全局截面 $s\ne0$。此时 $0\ne f^*s\in H^0(U,\cO_U(D+f^*H))$，且 $F:=(f^*s=0)\sim D+f^*H$。
\end{proof}

\begin{remark}{6.29}
若定理 6.27 中的 $X$ 是代数簇，则证明中只需用到射影簇的 MMP。

首先，与上面的证明同样，选取到非奇异射影曲线 $V$ 的开集 $V^\circ$ 上的态射 $\pi:(P\in X)\to(0\in V^\circ)$。

其次，考察同时奇点消解（定理 4.28）的证明中定理 4.39 是如何使用的。对于 $(0\in V^\circ)$，适当选取概形的平展态射 $(0\in V^\#)\to(0\in V^\circ)$，则由命题 4.69 可诱导出从 $(0\in V^\#)$ 到 Du Val 奇点的形变空间 $U$ 的概形态射。在 $\C^n$ 的包含所关注点的开邻域上，存在具有同样性质的平展态射 $(\C^n)^\#\to\C^n$（定理 4.39）。因此可以找到非奇异曲线 $W$ 的有限覆盖 $g:W\to V$，使得 $W^\circ=g^{-1}(V^\circ)\to V^\circ$ 经由 $V^\circ\times_U(\C^n)^\#$ 分解。此时，$W^\circ$ 上存在两个 Du Val 奇点的代数形变族。其一是 $X\times_{V^\circ}W^\circ$，关于它需要证明 (1)。其二是 $\bar Y_n\times_{\C^n}W^\circ$，它是只具有 Du Val 奇点的射影曲面的射影形变族 $\bar Y_n\times_{\C^n}W$ 的开集。注意，$\bar Y_n\times_{\C^n}W$ 存在整体的射影同时极小奇点消解 $Y_n\times_{\C^n}W$。

这里，$X\times_{V^\circ}W^\circ$ 是否代数同构于 $\bar Y_n\times_{\C^n}W^\circ$ 尚不清楚，但在奇点附近解析同构则是已知的。（事实上，定理 4.28 的证明用到了这一事实。）

% [原文 p.247]
(1) 在 $X\times_{V^\circ}W^\circ$ 的情形证明，只需在其与 $\bar Y_n\times_{\C^n}W^\circ$ 局部解析同构的情形证明即可。然而后者存在射影的同时极小奇点消解。因此，若在证明中把 $h:Z\to X$ 换成所得的 $Y_n\times_{\C^n}W\to\bar Y_n\times_{\C^n}W$，就只需射影簇的 MMP。
\end{remark}

\begin{definition}{6.30}
设 $X$ 是至多具有标准奇点的三维簇。以 $e(X)$ 表示 $X$ 上的例外除子 $E$ 中满足 $a(E,X)=0$ 者的个数。注意由命题 2.36 有 $e(X)<\infty$。（但若 $X$ 是解析空间，则须把 $X$ 限制到紧致部分集合 $K$ 的充分小邻域上。）
\end{definition}

\begin{proof}[定理 6.23 的证明]
对 $e(X)$ 归纳。若 $e(X)=0$，注意 $X$ 是终端的。若有 $X$ 上的点 $P$ 使 $(P\in X)$ 的指数 1 覆盖不是 cDV 奇点，则由定理 5.35，存在射影双有理态射 $f:Y\to X$ 满足 $K_Y=f^*K_X$，且至少把一个除子收缩到 $P$。因此 $e(Y)<e(X)$，由归纳法此情形已完成。

其次，假定任意 $(P\in X)$ 的指数 1 覆盖都是 cDV 奇点，且进一步假定 $\dim\operatorname{Sing}(X)=1$。由定理 6.27，存在射影双有理态射 $f:Y\to X$ 满足 $K_Y=f^*K_X$，且至少把一个除子收缩到曲线。有 $e(Y)<e(X)$，此情形也完成。

最后剩下的是任意 $(P\in X)$ 的指数 1 覆盖都是 cDV 奇点、且 $\dim\operatorname{Sing}(X)=0$ 的情形。但此时由引理 6.31，$X$ 已经是终端的，需要证明的已全部完成。

这里指出，由定理 5.35，所给的构造具有注意 6.24 中所述的性质。

若 $X$ 是解析的，则在紧致部分集合 $K$ 的充分小邻域上讨论，之后再变动 $K$。由注意 6.24，这些构造自然相容，从而可知所得簇在 $X$ 的任意紧致部分集合的适当邻域上是射影的。
\end{proof}

% [原文 p.248]
\begin{lemmap}[{[Rei83b]}]{6.31}
设 $\pi:(\tilde P\in\tilde X)\to(P\in X)$ 是标准奇点的任意次指数 1 覆盖（定义 5.19）。若 $(\tilde P\in\tilde X)$ 是终端的，则 $(P\in X)$ 也是终端的。
\end{lemmap}

\begin{proof}
设 $\omega$ 是 $\cO_X(rK_X)$ 的生成元。此时 $\pi$ 是取 $\omega$ 的 $r$ 次根所得的循环覆盖（定义 5.19）。设 $g:W\to X$ 是奇点消解，$E$ 是例外除子，$w$ 是 $E$ 的泛点。设 $\Omega$ 是 $w$ 处 $\cO_W(K_W)$ 的局部生成元。设 $\sigma:\tilde W\to W$ 是 $W$ 在 $\tilde X$ 的函数域中的正规化。于是得到如下交换图。
\[
\begin{tikzcd}
\tilde W \arrow[r,"\sigma"] \arrow[d,"\tilde g"] & W \arrow[d,"g"] \\
\tilde X \arrow[r,"\pi"] & X.
\end{tikzcd}
\]
这里假定 $a(E,X)=0$。这意味着可用在 $w$ 的邻域中不为零的函数 $u$ 把 $g^*\omega$ 表为 $g^*\omega=u\cdot\Omega^{\otimes r}$。由于 $\sigma$ 是取 $u$ 的 $r$ 次根（并进一步正规化）所得的操作，故可知 $\sigma$ 在 $w$ 上的某点 $\tilde w$ 处是平展的。设 $\tilde E\subset\sigma^{-1}(E)$ 是包含 $\tilde w$ 的素除子。以上可知 $rK_{\tilde W}=\sigma^*g^*rK_X=\tilde g^*rK_{\tilde X}$ 在 $\tilde w$ 的邻域成立，且 $a(\tilde E,\tilde X)=0$。这与 $\tilde X$ 是终端的假设矛盾，故证明完成。
\end{proof}

定理 6.25 的证明需要如下两个结果。

\begin{theorem*}[{[Kaw88]}]{6.32}
设 $X$ 是至多具有终端奇点的三维代数簇（或解析空间），$D$ 是 Weil 除子。此时成立：
\begin{enumerate}[label=（\arabic*）,leftmargin=3.2em]
\item $R(X,D)=\bigoplus_{\nu\ge0}\cO_X(\nu D)$ 是有限生成的 $\cO_X$-代数层。
\item 存在射影（或固有）双有理态射 $g:Y\to X$ 满足如下条件：$Y$ 正规，$\operatorname{Ex}(g)$ 的余维至少 2，且 $D':=g_*^{-1}D$ 是 Q-Cartier 的并在 $X$ 上（或在 $X$ 的任意紧致部分集合的适当邻域上）是 $g$-丰富的。
\end{enumerate}
\end{theorem*}

% [Tutor Note: 定理 6.32 中的 $g$ 称为符号爆破（symbolic blowup），它把“符号幂代数的有限生成性”这一代数条件几何化为一个除子 $D'=g_*^{-1}D$；与终端化一样，它对局部化与解析化都保持相容。]

\begin{remark}{6.33}
定理 6.32 所给出的 $g:Y\to X$（称为由 $X$ 的 $-D$ 决定的符号爆破（symbolic blowup））具有与注意 6.24 中所述基准终端化同样的性质，这是显然的。
\end{remark}
